\subsection{A centerless rank-five Morita partner excludes \texorpdfstring{$R_{8,3,439}$}{R8,3,439}}
\label{subsec:r439-centerless-morita}

The preceding pointed-associator calculation makes full $S_3$
condensation unconditional for $R_{8,3,439}$.  The complete dual
based-ring enumeration and frozen rank-five exclusions leave only
$R_{5,3,40}$.  Its basis $(1,s,V,P,Q)$ is self-dual, its dimensions are
$(1,1,2,3+2\sqrt3,3+2\sqrt3)$, and its dimension is
$D=48+24\sqrt3$.  The Galois reduction permits us to take the original
category pseudo-unitary; its Morita dual is then pseudo-unitary by
Morita invariance of categorical and Frobenius--Perron dimensions
\cite[Theorem~2.15 and Corollary~8.14]{ENO};
its canonical spherical structure is provided by
\cite[Proposition~8.23]{ENO}.
Its center has no nontrivial invertibles.  The integral-center
criterion therefore forbids every nonunit integral-dimensional
central simple.

The formal codegrees of this dual ring are $D$, $D'=48-24\sqrt3$,
$3$, $4$, $4$.  The complete induced-unit search gives $24$ low
frames, of which five pass the ordinary moment conditions and just
one passes full distinguished $S$ symmetry:
\[
 (1,0,0,0,0),\quad(1,0,0,1,1),\quad(1,1,1,2,2),\quad
 (1,0,1,1,2),\quad(1,0,1,2,1).
\]
The two assignments of the equal-codegree characters both fix the
simple $A$ corresponding to $D'$.  The remaining Gram matrix is
\[
 Q=\begin{pmatrix}4&2&4&4\\2&6&7&7\\4&7&34&31\\4&7&31&34\end{pmatrix}.
\]
The exact high domain uses positivity of $Q-vv^{\mathsf T}$, the
central simple dimension $d$-number condition
\cite[Theorem~1.8]{OstrikCode}, and algebraic integrality of
$D/d_Z$ \cite[Proposition~3.3]{ENO}.
There are $38$ possible high restriction vectors.  Four have integral
dimension and one violates algebraic integrality of a distinguished
$S$-ratio.  Two independent complete integer enumeration algorithms
find exactly $58$ multiplicity patterns on the remaining vectors.

Let $\chi=(1,1,2,3-2\sqrt3,3-2\sqrt3)$ be the conjugate-FP character.
The distinguished simple $A$ is self-dual: it is the unique simple
of dimension $D/D'=7+4\sqrt3$ among the constituents of the
self-dual object $I(1)$ \cite[Theorem~2.13]{OstrikPiv3}.
The distinguished-row formula \cite[equation~(39)]{DeshpandeCenter}
and the Verlinde formula give
\[
 N_{A,A}^{A}=\frac{D}{(D')^3}
       \sum_{Z\in\operatorname{Irr} Z(\mathcal C)}
       \frac{\chi(FZ)^3}{\operatorname{FPdim}(Z)}.
\]
Counting the five low simples once and each high restriction with its
integer multiplicity gives the following complete list:
\[
\begin{array}{c|ccccc}
N_{A,A}^{A}&(7-\sqrt3)/6&1/3&(19-5\sqrt3)/12&(6-\sqrt3)/3&(1+2\sqrt3)/3\\
\text{number of patterns}&16&16&16&6&4
\end{array}
\]
No entry is an integer.  Thus no possible associator on the forced
rank-five dual can have the required center.  Consequently
$R_{8,3,439}$ admits no complex fusion-category realization.

The exact replay reconstructs the complete low-frame domain, both
integer enumerations, and every cubic value; a Fibonacci-center
control checks the cubic normalization.  Its canonical entry point is
\begin{quote}\small\path{evidence/rank8_frontier_20261021/r439_centerless_rank5/verlinde_exclusion.py}\end{quote}
The earlier rational height-six rotation and $STS$ relaxation has a
feasible branch; the integer cubic test supplies the additional
obstruction needed here.
